\documentclass[journal,10pt,compsoc,onecolumn]{IEEEtran}
\usepackage[final]{graphicx}
\usepackage{booktabs}
\usepackage{subcaption}
\usepackage{geometry}
\usepackage{amsmath}
\usepackage{hyperref}

\geometry{letterpaper, margin=1in}

% Setup graphics path. Put all your PNGs into a folder named "figures" in Overleaf
\graphicspath{{figures/}}

\begin{document}

\title{Empirical Evaluation of Multi-Agent Routing Topologies and Workflows for Legal Compliance Analysis}

\author{Charbel~\dots \\
\textit{Legal AI Lab}}

\maketitle

\begin{abstract}
As large language models (LLMs) are increasingly deployed in highly regulated domains like law and regulatory compliance, the choice of agentic orchestration architecture becomes a critical design parameter. This paper presents a comprehensive empirical evaluation of six distinct multi-agent routing topologies (Sequential, Routed Single-Agent, Parallel Execution, Legal-News Parallel, Legal-First Selective, and Verification Bypass) evaluated on compliance checks for the EU AI Act. Using a suite of 32 operational, text quality (NLG), and structural metrics, we analyze the speed-quality trade-offs inherent in these topologies. Our results indicate that routed single-agent systems maintain a high balance of information retrieval precision and execution speed, while sequential workflows exhibit high latency overhead without commensurate quality improvements.
\end{abstract}

\section{Introduction}
Modern AI orchestration frameworks (such as LangGraph) enable complex state-machine routing between specialized agents. However, quantitative comparisons of different routing Topologies are scarce. This paper provides a rigorous analysis of the performance, quality, and structural features of outputs produced by various multi-agent routing topologies.

\section{Routing Topologies \& Methodology}
We evaluate six routing configurations:
\begin{enumerate}
    \item \textbf{ALL (Sequential)}: Executes all retrieval and analysis nodes in sequence.
    \item \textbf{SINGLE (Router)}: A routing node dynamically selects the single most appropriate expert agent.
    \item \textbf{PARALLEL}: Spawns all retrieval and logic agents concurrently.
    \item \textbf{LEGAL-NEWS PAR.}: Runs the legal and news retrieval agents in parallel, bypassing general Q\&A.
    \item \textbf{LEGAL-FIRST}: Executes the legal agent first, then decides if news or general QA is required.
    \item \textbf{VERIFY-ONLY (Bypass)}: Routes directly to aggregation and verification, bypassing retrieval nodes.
\end{enumerate}

Outputs are evaluated against an expert-curated "Gold Standard" compliance check for the EU AI Act using a multi-dimensional metric framework comprising 32 metrics.

\section{Evaluation Results}
The computed metrics across all configurations are summarized in Table~\ref{tab:agent_workflows_eval}.

\begin{table*}[htbp]
\centering
\caption{Comparative Analysis of Multi-Agent Routing Topologies and Workflows}
\label{tab:agent_workflows_eval}
\small
\begin{tabular}{lcccccc}
\toprule
\textbf{Metric} & \textbf{ALL (Sequential)} & \textbf{SINGLE (Router)} & \textbf{PARALLEL} & \textbf{LEGAL-NEWS PAR.} & \textbf{LEGAL-FIRST} & \textbf{VERIFY-ONLY (Bypass)} \\
\midrule
\multicolumn{7}{l}{\textit{Operational \& Latency Metrics}} \\
\textbf{End-to-End Latency (s)} & 79.07 & 46.51 & 71.00 & 69.49 & 83.41 & 22.27 \\
\textbf{Backend Graph Latency (s)} & 67.52 & 45.44 & 69.85 & 68.40 & 82.27 & 20.82 \\
\textbf{Graph Execution Steps} & 20 & 16 & 20 & 18 & 18 & 14 \\
\textbf{Mean Step Latency (ms)} & 3376.2 & 2839.8 & 3492.5 & 3799.9 & 4570.5 & 1487.4 \\
\textbf{Generation Speed (words/s)} & 1.49 & 3.48 & 2.15 & 1.84 & 1.61 & 3.64 \\
\textbf{Estimated Tokens (Words * 1.3)} & 153 & 210 & 198 & 166 & 174 & 105 \\
\textbf{Simulated Run Cost (USD)} & \$0.0001 & \$0.0001 & \$0.0001 & \$0.0001 & \$0.0001 & \$0.0001 \\
\midrule
\multicolumn{7}{l}{\textit{Text Quality \& NLG Reference-Matching (vs. Gold Standard)}} \\
\textbf{BLEU-1 Score} & 0.1206 & 0.2218 & 0.1940 & 0.1473 & 0.1557 & 0.0441 \\
\textbf{BLEU-4 Score} & 0.0175 & 0.0399 & 0.0285 & 0.0257 & 0.0168 & 0.0066 \\
\textbf{ROUGE-1 F1-Score} & 0.3030 & 0.3570 & 0.3302 & 0.3259 & 0.3276 & 0.2611 \\
\textbf{ROUGE-2 F1-Score} & 0.0914 & 0.0782 & 0.0748 & 0.1042 & 0.0639 & 0.0670 \\
\textbf{ROUGE-L F1-Score} & 0.1667 & 0.1739 & 0.1628 & 0.1877 & 0.1565 & 0.1500 \\
\textbf{Word Jaccard Similarity} & 0.1850 & 0.2336 & 0.1728 & 0.1940 & 0.1891 & 0.1564 \\
\textbf{Character Jaccard Similarity} & 0.5532 & 0.7551 & 0.7174 & 0.6957 & 0.6596 & 0.6087 \\
\textbf{TF Cosine Similarity} & 0.4958 & 0.4981 & 0.6402 & 0.5683 & 0.5734 & 0.5743 \\
\textbf{Levenshtein Similarity} & 0.2520 & 0.2706 & 0.2716 & 0.2783 & 0.2582 & 0.1924 \\
\midrule
\multicolumn{7}{l}{\textit{Text Structural \& Domain-Specific Metrics}} \\
\textbf{Word Count (Verbosity)} & 118 & 162 & 153 & 128 & 134 & 81 \\
\textbf{Sentence Count} & 8 & 10 & 8 & 9 & 10 & 4 \\
\textbf{Paragraph Count} & 4 & 5 & 3 & 3 & 4 & 4 \\
\textbf{Lexical Diversity (TTR)} & 0.7456 & 0.7548 & 0.6824 & 0.7561 & 0.7795 & 0.7692 \\
\textbf{Temporal References Count} & 0 & 4 & 5 & 2 & 2 & 0 \\
\textbf{Bullet Points Count} & 4 & 6 & 5 & 5 & 7 & 3 \\
\bottomrule
\end{tabular}
\end{table*}

\subsection{Latency and Operational Profiles}
As shown in Fig.~\ref{fig:macro_latency}, the execution latency varies heavily. While the verification-only workflow completes in only 22.27 seconds, it lacks comprehensive retrieval context. The sequential workflow (ALL) takes 79.07 seconds due to its step-by-step nature. Figure~\ref{fig:graph_efficiency} displays the step execution, cost, token, and generation speed profiles. Figure~\ref{fig:stacked_latencies} presents the detailed stacked breakdowns of individual node execution across workflows.

\begin{figure*}[htbp]
  \centering
  \begin{subfigure}[b]{0.48\textwidth}
    \centering
    \includegraphics[width=\textwidth]{chart1_e2e_latency.png}
    \caption{End-to-End Latency}
    \label{fig:e2e_latency}
  \end{subfigure}
  \hfill
  \begin{subfigure}[b]{0.48\textwidth}
    \centering
    \includegraphics[width=\textwidth]{chart2_backend_latency.png}
    \caption{Backend Graph Latency}
    \label{fig:backend_latency}
  \end{subfigure}
  \\ \vspace{1em}
  \begin{subfigure}[b]{0.48\textwidth}
    \centering
    \includegraphics[width=\textwidth]{chart3_network_overhead.png}
    \caption{Network \& API Overhead}
    \label{fig:net_overhead}
  \end{subfigure}
  \hfill
  \begin{subfigure}[b]{0.48\textwidth}
    \centering
    \includegraphics[width=\textwidth]{chart5_latency_per_step.png}
    \caption{Mean Latency per Step}
    \label{fig:latency_per_step}
  \end{subfigure}
  \caption{Macro Latency and Execution Profiles across Configurations.}
  \label{fig:macro_latency}
\end{figure*}

\begin{figure*}[htbp]
  \centering
  \begin{subfigure}[b]{0.48\textwidth}
    \centering
    \includegraphics[width=\textwidth]{chart4_execution_steps.png}
    \caption{State Transitions (Steps)}
  \end{subfigure}
  \hfill
  \begin{subfigure}[b]{0.48\textwidth}
    \centering
    \includegraphics[width=\textwidth]{chart6_generation_speed.png}
    \caption{Generation Speed (Words/s)}
  \end{subfigure}
  \\ \vspace{1em}
  \begin{subfigure}[b]{0.48\textwidth}
    \centering
    \includegraphics[width=\textwidth]{chart7_token_usage.png}
    \caption{Estimated Generated Tokens}
  \end{subfigure}
  \hfill
  \begin{subfigure}[b]{0.48\textwidth}
    \centering
    \includegraphics[width=\textwidth]{chart8_run_cost.png}
    \caption{Simulated Operational Cost}
  \end{subfigure}
  \caption{Graph Efficiency and Operational Metrics.}
  \label{fig:graph_efficiency}
\end{figure*}

\begin{figure*}[htbp]
  \centering
  \includegraphics[width=0.85\textwidth]{chart9_node_latency_stacked.png}
  \caption{Cumulative Stacked Bar Chart of Node-level Latency profiles across Workflows.}
  \label{fig:stacked_latencies}
\end{figure*}

\subsection{Text Quality \& Reference Alignment}
To measure how well output compliance guides align with the Gold Standard, we evaluate n-gram overlap. Figure~\ref{fig:nlg_quality} illustrates BLEU and ROUGE profiles.

\begin{figure*}[htbp]
  \centering
  \begin{subfigure}[b]{0.48\textwidth}
    \centering
    \includegraphics[width=\textwidth]{chart18_bleu_scores.png}
    \caption{BLEU-1 to BLEU-4 Profile}
  \end{subfigure}
  \hfill
  \begin{subfigure}[b]{0.48\textwidth}
    \centering
    \includegraphics[width=\textwidth]{chart19_rouge1_scores.png}
    \caption{ROUGE-1 Overlap}
  \end{subfigure}
  \\ \vspace{1em}
  \begin{subfigure}[b]{0.48\textwidth}
    \centering
    \includegraphics[width=\textwidth]{chart20_rouge2_scores.png}
    \caption{ROUGE-2 Overlap}
  \end{subfigure}
  \hfill
  \begin{subfigure}[b]{0.48\textwidth}
    \centering
    \includegraphics[width=\textwidth]{chart21_rougeL_scores.png}
    \caption{ROUGE-L Longest Common Subseq.}
  \end{subfigure}
  \caption{Text Quality and Overlap Comparison against Gold Standard.}
  \label{fig:nlg_quality}
\end{figure*}

\subsection{Trade-Off Profiles}
Figure~\ref{fig:multidim_analysis} illustrates the speed-quality trade-off space and provides a spider-radar comparison showing that the SINGLE workflow dominates across multiple dimensions.

\begin{figure*}[htbp]
  \centering
  \begin{subfigure}[b]{0.48\textwidth}
    \centering
    \includegraphics[width=\textwidth]{chart28_quality_vs_latency.png}
    \caption{Quality (ROUGE-L) vs. Latency Trade-off}
  \end{subfigure}
  \hfill
  \begin{subfigure}[b]{0.48\textwidth}
    \centering
    \includegraphics[width=\textwidth]{chart29_performance_radar.png}
    \caption{Multi-dimensional Performance Radar}
  \end{subfigure}
  \caption{Comprehensive Multi-dimensional Trade-off Profiles.}
  \label{fig:multidim_analysis}
\end{figure*}

\section{Conclusion}
This paper presented a detailed empirical evaluation of agent topologies. We conclude that parallel and routed single-agent workflows yield the highest response fidelity and alignment with references, while sequential setups suffer from prohibitive latency without measurable quality gains.

\end{document}
